\chapter*{Pref\'acio}
\addcontentsline{toc}{chapter}{Pref\'acio}

Este livro conta uma hist\'oria. N\~ao \'e uma cole\c{c}\~ao de
receitas isoladas: cada cap\'itulo introduz um personagem, uma tens\~ao
ou uma virada que s\'o fazem sentido luz do que veio antes. O personagem
central \'e um sistema de IA corporativo que nasce simples --- uma
chamada ao LLM --- e amadurece progressivamente at\'e operar em
produ\c{c}\~ao banc\'aria com rastreabilidade, avalia\c{c}\~ao cont\'inua
e governan\c{c}a t\'ecnica. O leitor que acompanhar essa jornada
cap\'itulo a cap\'itulo sair\'a capaz n\~ao apenas de reproduzir o
c\'odigo, mas de tomar as mesmas decis\~oes de projeto --- e de
justific\'a-las.

\section*{A hist\'oria em dois ciclos}

Toda plataforma de IA em produ\c{c}\~ao vive em dois ciclos
interdependentes, que este livro chama de \textbf{loop interno} e
\textbf{loop externo}.

O \textbf{loop interno} (\textit{inner loop}) \'e o ciclo de
\textit{desenvolvimento}: experimento, avalia\c{c}\~ao offline,
valida\c{c}\~ao humana e promo\c{c}\~ao para produ\c{c}\~ao. Ele \'e
onde o engenheiro vive --- no notebook, nos testes unit\'arios, na
constru\c{c}\~ao do \textit{golden dataset}, na calibra\c{c}\~ao do
juiz LLM. Os cap\'itulos 1 a 10 narram, passo a passo, a constru\c{c}\~ao
desse loop.

O \textbf{loop externo} (\textit{outer loop}) \'e o ciclo de
\textit{produ\c{c}\~ao}: a plataforma responde a clientes reais, mede
confiabilidade em tempo real, coleta evid\^encias e as devolve ao
loop interno para reajuste. Os cap\'itulos 11 a 15 instrumentam esse
ciclo --- observabilidade, m\'etricas de mercado, LLM-as-Judge,
\textit{golden dataset} e \textit{roadmap} de evolu\c{c}\~ao cont\'inua.
Os cap\'itulos 16 a 18 garantem que o ciclo se sustente com qualidade,
seguran\c{c}a e CI/CD.

Os dois loops n\~ao s\~ao seq\"uenciais: eles giram em paralelo e
se alimentam mutuamente. O cap\'itulo 10 \'e o ponto de converg\^encia
onde essa din\^amica se torna expl\'icita pela primeira vez.

\section*{Como ler este livro}

\begin{itemize}[leftmargin=*, noitemsep]
  \item \textbf{Leitura linear:} recomendada na primeira passagem.
        Cada cap\'itulo pressup\~oe os anteriores; o glossario e os
        pr\'e-requisitos de cada cap\'itulo indicam o m\'inimo
        necess\'ario.
  \item \textbf{Leitura por camada:} caps.\ 1--4 (fundamentos) $\to$
        caps.\ 5--9 (plataforma) $\to$ caps.\ 10--14 (qualidade) $\to$
        caps.\ 15--18 (opera\c{c}\~ao). Cada camada pode ser retomada
        independentemente ap\'os a leitura completa.
  \item \textbf{Rastreabilidade:} cada cap\'itulo mant\'em a tríade
        \textit{conceito $\to$ implementa\c{c}\~ao $\to$ evidencia de
        qualidade}. Os planos de teste e os crit\'erios de aceite s\~ao
        parte do texto, n\~ao ap\^endice.
\end{itemize}

\section*{\'Indice proposto}

\noindent\textbf{Loop interno --- constru\c{c}\~ao} (caps.\ 1--10)
\begin{enumerate}[leftmargin=*]
  \item Fundamentos de IA e LLMs
  \item Ambiente, ferramentas e reprodutibilidade
  \item Agentes com LangGraph e mem\'oria
  \item Tools, skills e MCP
  \item Fundamentos de RAG
  \item RAG corporativo e arquitetura de dados
  \item Pipeline ass\'incrono em AWS
  \item Segmenta\c{c}\~ao, personas e digital twins
  \item GraphRAG com Neptune
  \item \textbf{Pipeline integrado em produ\c{c}\~ao} $\leftarrow$ \textit{ponto de converg\^encia}
\end{enumerate}

\noindent\textbf{Loop externo --- qualidade e opera\c{c}\~ao} (caps.\ 11--18)
\begin{enumerate}[leftmargin=*, start=11]
  \item Confiabilidade, observabilidade e opera\c{c}\~ao
  \item M\'etricas de avalia\c{c}\~ao de mercado
  \item LLM as Judge
  \item Golden dataset e regress\~ao de qualidade
  \item Roadmap de produ\c{c}\~ao e evolu\c{c}\~ao cont\'inua
  \item CI/CD multiambiente banc\'ario
  \item Qualidade e seguran\c{c}a nas esteiras
  \item Frontend, CloudFront e console de opera\c{c}\~ao
\end{enumerate}
